% BHU PYQ -- Professional Exam Template
\documentclass[11pt,a4paper]{article}

\usepackage[a4paper,top=2.5cm,bottom=2.5cm,left=2.8cm,right=2.8cm,headheight=14pt]{geometry}
\usepackage{amsmath,amssymb}
\usepackage[T1]{fontenc}
\usepackage[utf8]{inputenc}
\usepackage{lmodern}
\usepackage{microtype}
\usepackage{fancyhdr}
\usepackage{enumitem}
\usepackage{hyperref}
\usepackage{lastpage}
\usepackage{parskip}

\hypersetup{colorlinks=true,urlcolor=black,linkcolor=black,
  pdftitle={B.Sc. (Hons.) Botany Sem-V BOB-501 -- BHU PYQ}}

\pagestyle{fancy}
\fancyhf{}
\renewcommand{\headrulewidth}{0.4pt}
\renewcommand{\footrulewidth}{0.4pt}
\fancyhead[L]{\small   \textit{Banaras Hindu University}}
\fancyhead[R]{\small   \textit{Botany Paper: BOB-501}}
\fancyfoot[C]{\small Page  \thepage\ of \pageref{LastPage}}


\newlist{parts}{enumerate}{2}
\setlist[parts,1]{label=(\alph*),leftmargin=*,itemsep=4pt,topsep=3pt}
\setlist[parts,2]{label=(\roman*),leftmargin=*,itemsep=2pt,topsep=2pt}

\newcommand{\pts}[1]{\hfill   \textit{[#1]}}


\begin{document}


\begin{center}
  {\large\bfseries BANARAS HINDU UNIVERSITY}\\[2pt]
  {\normalsize B.Sc. (Hons.) Semester V Examination, 2016-17}\\[6pt]
  \rule{0.8\linewidth}{0.5pt}\\[6pt]
  {\large\bfseries Botany}\\[4pt]
  {\large\bfseries Paper: BOB-501 --- Comparative Studies of Cryptogams}\\[4pt]
  {\normalsize Time: Three Hours \hfill Full Marks: 70}
\end{center}

\rule{\linewidth}{0.4pt}

\smallskip



\noindent   \textit{Note: Attempt five questions in all, selecting one question each from Section-A, B, C, and D. Section-E is compulsory. The figures in the right-hand margin indicate marks.}

\bigskip

\begin{center}
   \textbf{SECTION -- A}
\end{center}

\medskip


\noindent   \textbf{Question 1.} Give an illustrated account of thallus organization in Chlorophyceae.\pts{15}

\medskip


\noindent   \textbf{Question 2.} Write explanatory notes on any two of the following:\pts{$7.5  \times 2 = 15$}

\begin{parts}
  \item Triphasic life history in a red alga
  \item Photosynthetic pigments in major classes of algae
  \item Oogamous reproduction in any brown alga
\end{parts}

\bigskip

\begin{center}
   \textbf{SECTION -- B}
\end{center}

\medskip


\noindent   \textbf{Question 3.} With the help of suitable diagrams, give a critical account of types of fungal spores and modes of their liberation.\pts{15}

\medskip


\noindent   \textbf{Question 4.} Write notes on any two of the following:\pts{$7.5  \times 2 = 15$}

\begin{parts}
  \item Classification of Ascomycotina
  \item Major types of life cycles in fungi
  \item Importance of fungi in industry and agriculture
\end{parts}

\bigskip

\begin{center}
   \textbf{SECTION -- C}
\end{center}

\medskip


\noindent   \textbf{Question 5.} Give an illustrated account of evolution of sporophyte in Bryophyta by progressive sterilization of sporogenous tissues.\pts{15}

\medskip


\noindent   \textbf{Question 6.} Write notes on any two of the following:\pts{$7.5  \times 2 = 15$}

\begin{parts}
  \item Vegetative propagation in Bryophyta
  \item Structure and anatomy of gametophytes of   \textit{Pellia} and   \textit{Notothylas}
  \item Life history of   \textit{Sphagnum}
\end{parts}

\bigskip

\begin{center}
   \textbf{SECTION -- D}
\end{center}

\medskip


\noindent   \textbf{Question 7.} Give a critical account of structure of   \textit{Marsilea} sporocarp.\pts{15}

\medskip


\noindent   \textbf{Question 8.} Write notes on any two of the following:\pts{$7.5  \times 2 = 15$}

\begin{parts}
  \item Major types of steles in Pteridophyta
  \item Apogamy and apospory in ferns
  \item Salient features of   \textit{Psilotum}
\end{parts}

\pagebreak

\begin{center}
   \textbf{SECTION -- E}
\end{center}

\medskip


\noindent   \textbf{Question 9.} Answer the following in not more than 50 words each:\pts{$1  \times 10 = 10$}

\begin{parts}
  \item Name two important food reserve products found in Cyanophyceae.
  \item Give one example each of algae showing isomorphic and heteromorphic alternation of generations.
  \item Give examples of uniaxial and multiaxial thallus in algae.
  \item What is diplanetism?
  \item Name different types of spores produced by   \textit{Puccinia}.
  \item Name a bryophyte which has elaterophores in the sporogonium.
  \item What are retort cells?
  \item Write down the name of a ligulate and heterosporous pteridophyte.
  \item What is the function of peristome teeth?
  \item What is false indusium?
\end{parts}

\vfill
\rule{\linewidth}{0.4pt}

\begin{center}\small
\href{https://bhu-exam.vercel.app/}{Click Here}\\[4pt]\href{https://me-aryan.vercel.app/}{Click Here}\\[4pt]\href{https://quashan-me.vercel.app/}{Click Here}
\end{center}

\end{document}
